\cotefile{148}
\scanpage{3}

One way of seeing it is to say that
the datum of $\uline{s}$ defines a splitting
\inserted{$\varphi_{\uline{s}}$} of
the extension $(*)$ above, \unsure{and that viewed as|and as}
subgroup of $\operatorname{SRT}(X)$,
$S_1T(X,\uline{s})$ is identified
with the \unsure{inverse|\deleted{of }inverse} image of
\[
\varphi_{\uline{s}}(T(X))\subset\operatorname{SUT}(X)
\]
under the canonical homomorphism
$\operatorname{SUR}(X)\mathrel{\unsure{\text{---}|\longrightarrow}}\operatorname{SUT}(X)$
of passage to the quotient. The change of
$\uline{s}$ to $\uline{s}'=u.\uline{s}$
transforms $\varphi_{\uline{s}}$ into
\[
\varphi_{\uline{s}'}=\operatorname{int}(u)\circ\varphi_{\uline{s}}
\]
hence (if $\widetilde u\in\mathbb R^I$ lifts $u$\unsure{,|)} one finds
\[
S_1T(X,\uline{s}')
=\operatorname{int}(\widetilde u)S_1T(X,\uline{s})\unsure{.|,}
\]

In terms of $\unsure{\overline X|\vec X}$, the definition of
$S_1T(X,\uline{s})$
depends on a \uline{non-innocuous} choice $\uline{s}$,
on the other hand that of $S_1RT(X)$ depends only on
an innocuous choice, that of the family of
toric structures $(t_i)_{i\in I}$.
Moreover $\prod T_i$ as
a torsor under $\unsure{\mathbb U^I|\mathbb M^I}$ is realized
\inserted{(not faithfully, incidentally)}
as \unsure{a class|one of the} of splittings of the extension $(*)$
(via inner automorphisms by $\mathbb U^I$)
$\bigl[$deduced from $\operatorname{SRT}(X)$
with the additional structures$\unsure{\bigr\}|\bigr]}$
namely\unsure{)|,} the inclusions

\diagram{
\begin{tikzpicture}[baseline=(current bounding box.center)]
\node (a) at (0,1.65) {$\mathbb{R}^{I}$};
\node (b) at (0,-1.65) {$S_1T^{!}(X)$};
\node (c) at (4.6,0) {$S\unsure{R|K}T(X)$};
\draw[->] (a.east) to[out=-35,in=165] (c.north west);
\draw (3.12,.73) .. controls (3.35,.66) and (3.52,.43) .. (3.73,.24);
\draw[->] (b.north east)--(c.south west);
\end{tikzpicture}
}

\unsure{coinciding|coinciding|coincide} on $\mathbb Z^I$.

\margin{2}
